\chapter[TOGAF-lite, ADM, dan Tinjauan Pustaka]{TOGAF-lite, ADM, dan Tinjauan Pustaka: Menyusun Visi Arsitektur dan Fondasi Akademik}
\label{ch:togaf-lite}
\setlength{\emergencystretch}{2em}

\section*{Tujuan Pembelajaran}
\begin{itemize}
  \item Menjelaskan fungsi TOGAF-lite dan ADM sebagai kerangka kerja ringan
        untuk menata perubahan arsitektur enterprise.
  \item Menyusun visi arsitektur awal yang berangkat dari masalah, pertanyaan
        riset, pemangku kepentingan, dan batasan organisasi.
  \item Merancang strategi pencarian pustaka dan bibliografi awal yang relevan
        dengan topik makalah individual.
  \item Menghubungkan sintesis pustaka, visi arsitektur, dan metode
        TOGAF-lite ke dalam rancangan awal makalah publikasi.
\end{itemize}

Bab ini terutama mendukung capaian pembelajaran yang berkaitan dengan
pemahaman konsep arsitektur enterprise, analisis pemangku kepentingan,
penyusunan argumen berbasis pustaka, dan penyiapan naskah publikasi
individual.

\section{Ringkasan Awal Bab}

Bab sebelumnya membantu pembaca memilih topik, merumuskan masalah, dan
membentuk pertanyaan riset awal. Bab ini mengubah rumusan tersebut menjadi
dua fondasi yang dapat dikerjakan secara terstruktur: visi arsitektur dan
tinjauan pustaka awal. Visi arsitektur menjelaskan arah perubahan yang ingin
dibahas dalam makalah. Tinjauan pustaka memberi dasar akademik agar arah
perubahan tersebut tidak hanya menjadi opini pribadi, melainkan berdiri di
atas konsep, temuan, standar, dan pengalaman yang sudah pernah dibahas dalam
literatur.

TOGAF digunakan secara ringan, bukan sebagai prosedur birokratis penuh. Yang
diambil adalah cara berpikir ADM: memahami konteks, merumuskan visi,
memetakan kondisi saat ini, membayangkan kondisi target, menemukan
kesenjangan, lalu menyusun arah perubahan. Dalam konteks mata kuliah ini, ADM
menjadi alat bantu untuk menulis makalah yang lebih rapi: masalahnya jelas,
kerangka konseptualnya terlihat, metode analisisnya dapat dijelaskan, dan
rekomendasinya dapat ditelusuri kembali ke bukti.

Pada akhir bab, setiap mahasiswa menghasilkan visi arsitektur satu halaman,
strategi pencarian pustaka, bibliografi awal minimal sepuluh rujukan, dan
kerangka metode singkat untuk makalah individual.

\section{Pendahuluan}

Makalah publikasi membutuhkan dua hal yang sering tertukar. Pertama, ia perlu
arah perubahan yang jelas. Kedua, ia perlu dasar pustaka yang cukup untuk
menunjukkan bahwa penulis memahami percakapan ilmiah dan profesional di
bidangnya. Jika hanya ada arah perubahan, naskah mudah berubah menjadi
proposal internal. Jika hanya ada pustaka, naskah mudah menjadi rangkuman
artikel tanpa kontribusi.

Bab ini mempertemukan keduanya. TOGAF-lite dipakai untuk menyusun arah
perubahan melalui visi arsitektur. Tinjauan pustaka dipakai untuk membangun
argumen, memilih konsep, dan menjelaskan mengapa usulan perubahan tersebut
layak dibahas. Dengan begitu, mahasiswa tidak sekadar menulis bahwa organisasi
perlu ``bertransformasi digital'', tetapi dapat menunjukkan perubahan apa yang
dibahas, bukti apa yang mendukung, dan bagaimana analisis dilakukan.

Pendekatan ini sesuai dengan keluaran utama mata kuliah: satu makalah
individual yang siap dikembangkan untuk jurnal atau konferensi. Artefak
arsitektur tetap dibuat, tetapi posisinya sebagai penguat naskah. Artefak
yang baik membantu penulis berpikir dan membantu pembaca memahami argumen.

\section{Skenario Pembuka dan Pertanyaan Pemandu}

Tim pengembangan layanan digital Bank Wanua mulai menyadari bahwa persoalan
mereka bukan sekadar aplikasi yang lambat atau data nasabah yang tersebar.
Masalah yang lebih besar adalah tidak adanya gambaran bersama tentang arah
perubahan. Unit bisnis ingin layanan baru cepat diluncurkan. Unit kepatuhan
menuntut tata kelola data yang lebih kuat. Tim teknologi ingin mengurangi
ketergantungan pada sistem lama. Pimpinan meminta satu dokumen ringkas yang
menjelaskan perubahan apa yang akan dilakukan, mengapa perubahan itu penting,
dan dasar akademik apa yang mendukung pilihan tersebut.

Seorang mahasiswa yang memilih kasus Bank Wanua sebagai topik makalah
menghadapi pertanyaan serupa. Ia sudah memiliki isu awal, misalnya integrasi
layanan digital, tata kelola data, atau modernisasi API. Namun, ia belum tahu
bagaimana menjadikan isu itu sebagai makalah yang layak dipublikasikan. Jika
hanya menulis deskripsi masalah, naskah akan terasa seperti laporan internal.
Jika hanya mengutip banyak artikel, naskah akan kehilangan arah praktis. Bab
ini membantu menautkan keduanya: arah perubahan dirumuskan melalui visi
arsitektur, sedangkan dasar akademiknya dibangun melalui tinjauan pustaka.

\subsection*{Pertanyaan Pemandu}

\begin{enumerate}
  \item Perubahan arsitektur apa yang benar-benar ingin dijelaskan dalam
        makalah?
  \item Siapa pemangku kepentingan utama, dan kekhawatiran apa yang perlu
        dijawab?
  \item Bagaimana ADM dapat disederhanakan agar cocok untuk analisis akademik
        individual?
  \item Literatur apa yang diperlukan agar argumen makalah tidak berdiri di
        atas asumsi pribadi?
  \item Bagaimana visi arsitektur dan tinjauan pustaka dapat menyatu dalam
        metode penulisan makalah?
\end{enumerate}

\section{Mengapa Bab Ini Penting}

Mahasiswa non-TI sering kali sudah memahami masalah organisasi dengan baik,
tetapi belum tentu terbiasa mengubah masalah tersebut menjadi argumen akademik
yang sistematis. Sebaliknya, mahasiswa juga dapat mengumpulkan banyak artikel
tanpa benar-benar tahu bagaimana artikel itu akan dipakai dalam analisis.
Keduanya berisiko menghasilkan naskah yang dangkal: kaya kutipan tetapi miskin
arah, atau kaya pengalaman tetapi lemah dasar ilmiah.

TOGAF-lite membantu mahasiswa menata arah perubahan. Tinjauan pustaka
membantu mahasiswa menata dasar pengetahuan. Keduanya perlu berjalan bersama.
Visi arsitektur tanpa pustaka dapat berubah menjadi klaim normatif. Tinjauan
pustaka tanpa visi arsitektur dapat berubah menjadi daftar ringkasan artikel.
Makalah publikasi membutuhkan keduanya: fokus perubahan yang tajam dan argumen
yang dapat dipertanggungjawabkan.

\section{Konsep Inti dan Glosarium Mini}

\subsection{TOGAF-lite}

TOGAF adalah kerangka kerja arsitektur enterprise yang menyediakan konsep,
siklus kerja, dan artefak untuk mengelola perubahan arsitektur organisasi
\cite{togaf10}. Dalam praktik profesional, TOGAF dapat digunakan secara luas
dan formal. Dalam mata kuliah ini, TOGAF digunakan secara ringan. Istilah
``TOGAF-lite'' berarti mahasiswa tidak diminta meniru seluruh dokumentasi
TOGAF, melainkan mengambil pola berpikir yang paling berguna untuk makalah:
memahami konteks, menyatakan visi, membandingkan kondisi saat ini dan kondisi
target, lalu menyusun rekomendasi perubahan.

Pendekatan ringan ini penting karena keluaran utama mata kuliah adalah makalah
individual, bukan proyek implementasi besar. Artefak arsitektur tetap dibuat,
tetapi posisinya sebagai alat bantu analisis. Artefak tersebut harus
memperjelas argumen naskah, bukan menjadi tumpukan lampiran yang terpisah dari
tulisan.

\subsection{ADM sebagai Siklus Berpikir}

Architecture Development Method atau ADM adalah bagian penting dari TOGAF. ADM
menggambarkan proses pengembangan arsitektur secara bertahap, mulai dari
persiapan, visi, arsitektur bisnis, sistem informasi, teknologi, peluang dan
solusi, perencanaan migrasi, implementasi, hingga tata kelola perubahan
\cite{togaf10}. Untuk kebutuhan makalah individual, ADM dapat disederhanakan
menjadi enam pertanyaan:

\begin{enumerate}
  \item Apa masalah dan pertanyaan risetnya?
  \item Apa visi arsitektur yang ingin dicapai?
  \item Bagaimana kondisi saat ini?
  \item Seperti apa kondisi target yang masuk akal?
  \item Kesenjangan apa yang perlu dijembatani?
  \item Rekomendasi perubahan apa yang paling layak dipertahankan secara
        akademik?
\end{enumerate}

\begin{figure}[htbp]
\centering
\scalebox{0.88}{
\begin{tikzpicture}[
  admphase/.style={rectangle, rounded corners=1.5mm, draw=praditagreen,
    fill=praditagreen!10, thick, align=center, text width=31mm,
    minimum height=12mm},
  flow/.style={-Latex, thick, draw=praditagreen},
  note/.style={rectangle, rounded corners=1.5mm, draw=gray!50,
    fill=gray!10, align=center, text width=33mm, minimum height=10mm}
]
  \node[admphase] (problem) {Masalah dan\\pertanyaan riset};
  \node[admphase, right=7mm of problem] (vision) {Visi\\arsitektur};
  \node[admphase, right=7mm of vision] (baseline) {Kondisi\\saat ini};
  \node[admphase, below=9mm of baseline] (target) {Kondisi\\target};
  \node[admphase, left=7mm of target] (gap) {Analisis\\kesenjangan};
  \node[admphase, left=7mm of gap] (roadmap) {Rekomendasi\\dan roadmap};

  \draw[flow] (problem) -- (vision);
  \draw[flow] (vision) -- (baseline);
  \draw[flow] (baseline) -- (target);
  \draw[flow] (target) -- (gap);
  \draw[flow] (gap) -- (roadmap);
  \draw[flow] (roadmap.west) .. controls +(-10mm,5mm) and +(-10mm,-5mm) ..
    (problem.west);

  \node[note, below=5mm of vision] (evidence) {Tinjauan pustaka\\dan bukti};
  \draw[flow, dashed] (evidence) -- (vision);
  \draw[flow, dashed] (evidence) -- (gap);
\end{tikzpicture}
}
\caption{Penyederhanaan ADM untuk analisis makalah individual.}
\label{fig:adm-lite-ch3}
\end{figure}

\subsection{Visi Arsitektur}

Visi arsitektur adalah pernyataan awal tentang arah perubahan yang ingin
dicapai. Visi ini tidak perlu panjang, tetapi harus cukup jelas untuk menjawab
tiga hal: masalah apa yang ingin diubah, hasil apa yang diharapkan, dan batas
mana yang tidak akan dibahas. Visi arsitektur juga perlu menyebut pemangku
kepentingan utama karena perubahan arsitektur hampir selalu memengaruhi lebih
dari satu pihak.

Dalam makalah, visi arsitektur membantu pembaca memahami posisi naskah sejak
awal. Jika topik mahasiswa adalah integrasi layanan kesehatan daerah, visi
arsitektur menjelaskan apakah fokusnya pada interoperabilitas data, alur
layanan pasien, tata kelola API, kepatuhan data pribadi, atau kombinasi yang
masih dapat dikelola. Tanpa visi semacam ini, makalah mudah melebar.

\subsection{Tinjauan Pustaka}

Tinjauan pustaka bukan kumpulan ringkasan artikel. Tinjauan pustaka adalah
cara menunjukkan bahwa penulis memahami percakapan akademik dan profesional
yang relevan dengan topiknya. Pustaka dipilih untuk menjawab kebutuhan
argumen: definisi konsep, temuan penelitian, kerangka kerja, praktik baik,
standar, regulasi, dan celah yang masih dapat dikaji.

Untuk bab ini, tinjauan pustaka awal sebaiknya memuat sedikitnya sepuluh
rujukan yang bermakna. Jumlah tersebut bukan angka magis, tetapi batas awal
agar mahasiswa tidak hanya bergantung pada satu atau dua sumber. Pada tahap
berikutnya, daftar ini dapat berkembang sesuai kebutuhan analisis dan target
jurnal atau konferensi.

\subsection{Manajer Referensi dan Bibliografi Kerja}

Manajer referensi membantu mahasiswa menyimpan metadata, mengelompokkan
sumber, menulis catatan ringkas, dan menghasilkan daftar pustaka secara
konsisten. Alat yang digunakan boleh berbeda, tetapi disiplin kerjanya sama:
setiap sumber yang disimpan harus memiliki alasan pemakaian. Bibliografi kerja
tidak cukup berisi judul dan tautan; ia perlu mencatat konsep utama, metode,
konteks, temuan penting, dan hubungan sumber tersebut dengan pertanyaan riset.

\subsection*{Glosarium Mini}

\begin{description}
  \item[ADM] Siklus berpikir dalam TOGAF untuk mengembangkan arsitektur
        enterprise secara bertahap.
  \item[TOGAF-lite] Pemakaian prinsip TOGAF secara selektif dan ringan untuk
        kebutuhan analisis makalah.
  \item[Visi arsitektur] Pernyataan arah perubahan arsitektur yang menjelaskan
        masalah, hasil yang diharapkan, ruang lingkup, dan pemangku
        kepentingan.
  \item[Concern] Kepentingan, kebutuhan, risiko, atau kekhawatiran pemangku
        kepentingan terhadap perubahan arsitektur.
  \item[Bibliografi kerja] Daftar pustaka yang masih berkembang dan disertai
        catatan kegunaan tiap sumber.
  \item[Matriks sintesis] Tabel yang membandingkan beberapa sumber berdasarkan
        konsep, metode, temuan, dan relevansinya dengan makalah.
\end{description}

\section{Kerangka Berpikir}

Bab ini menggunakan alur berpikir yang sederhana: pertanyaan riset menentukan
kebutuhan bukti; kebutuhan bukti mengarahkan pencarian pustaka; pustaka
membantu mempertajam visi arsitektur; visi arsitektur menjadi dasar analisis
kondisi saat ini, target, kesenjangan, dan rekomendasi.

\begin{figure}[htbp]
\centering
\scalebox{0.88}{
\begin{tikzpicture}[
  box/.style={rectangle, rounded corners=1.5mm, draw=praditagreen,
    fill=green!8, thick, align=center, text width=34mm, minimum height=12mm},
  link/.style={-Latex, thick, draw=praditagreen}
]
  \node[box] (rq) {Pertanyaan riset};
  \node[box, right=7mm of rq] (stake) {Pemangku kepentingan\\dan concern};
  \node[box, right=7mm of stake] (vision) {Visi arsitektur};
  \node[box, below=8mm of vision] (search) {Strategi pencarian\\pustaka};
  \node[box, left=7mm of search] (matrix) {Matriks sintesis};
  \node[box, left=7mm of matrix] (method) {Metode dan kerangka\\makalah};

  \draw[link] (rq) -- (stake);
  \draw[link] (stake) -- (vision);
  \draw[link] (vision) -- (search);
  \draw[link] (search) -- (matrix);
  \draw[link] (matrix) -- (method);
  \draw[link, dashed] (method) -- (rq);
\end{tikzpicture}
}
\caption{Hubungan antara visi arsitektur dan fondasi pustaka.}
\label{fig:vision-literature-loop}
\end{figure}

Kerangka ini sengaja dibuat berulang. Setelah membaca pustaka, mahasiswa
mungkin menyadari bahwa pertanyaan risetnya terlalu lebar. Setelah memetakan
pemangku kepentingan, mahasiswa mungkin melihat bahwa visi awalnya terlalu
teknis dan belum menjawab kebutuhan organisasi. Revisi semacam ini bukan tanda
kegagalan. Justru di situlah penulisan akademik mulai matang.

\section{Konteks Indonesia}

Makalah yang baik tidak hanya mengutip konsep global, tetapi juga memahami
konteks tempat masalah muncul. Untuk kasus Indonesia, konteks tersebut dapat
mencakup regulasi, kebijakan tata kelola data, karakter organisasi, tingkat
kematangan digital, serta ketersediaan sumber rujukan lokal. Misalnya, topik
tata kelola data pribadi perlu membaca kerangka perlindungan data pribadi
\cite{uupdp2022}. Topik sistem elektronik dan transaksi digital perlu
memperhatikan aturan penyelenggaraan sistem elektronik \cite{pp712019}.
Topik pertukaran data pemerintah dapat mengacu pada kerangka Satu Data
Indonesia \cite{satudata}. Topik layanan keuangan digital dapat
mempertimbangkan standar API pembayaran nasional \cite{snapbi}.

Pada saat yang sama, mahasiswa perlu berhati-hati. Dokumen regulasi bukan
pengganti artikel ilmiah. Regulasi memberi batas dan konteks; artikel ilmiah
membantu membangun konsep, metode, dan perbandingan. Keduanya saling
melengkapi. Untuk pencarian literatur, mahasiswa dapat mengombinasikan sumber
internasional, repositori nasional, artikel jurnal Indonesia, prosiding
konferensi, dokumen regulator, dan laporan institusi. Yang penting bukan
banyaknya sumber, melainkan kemampuan menjelaskan mengapa sumber itu dipakai.

\section{Studi Kasus Utama: Visi Arsitektur Bank Wanua}

Kasus Bank Wanua digunakan untuk menunjukkan bagaimana topik awal dapat diubah
menjadi visi arsitektur dan kebutuhan pustaka. Misalkan topik mahasiswa adalah
modernisasi layanan digital bank daerah. Rumusan masalah awalnya berbunyi:
``Layanan digital Bank Wanua sulit berkembang karena data nasabah, proses
persetujuan produk, dan kanal layanan masih tersebar di beberapa sistem
lama.''

Rumusan tersebut sudah berguna, tetapi belum cukup untuk makalah. Mahasiswa
perlu memperjelas arah perubahan. Salah satu visi arsitektur awal dapat
dirumuskan sebagai berikut:

\begin{quote}
Bank Wanua membutuhkan arsitektur layanan digital yang memungkinkan integrasi
data nasabah, orkestrasi proses produk, dan pembukaan API secara terkendali,
sehingga pengembangan layanan baru dapat berlangsung lebih cepat tanpa
mengabaikan keamanan, kepatuhan, dan keberlanjutan sistem lama.
\end{quote}

Visi ini memuat beberapa keputusan awal. Fokusnya bukan mengganti seluruh
sistem inti bank, melainkan mengelola integrasi, proses, dan API. Nilai
bisnisnya adalah percepatan layanan baru. Risikonya adalah keamanan,
kepatuhan, dan keberlanjutan sistem lama. Dari sini, kebutuhan pustaka menjadi
lebih jelas: mahasiswa perlu membaca tentang arsitektur enterprise, platform
digital, tata kelola API, modernisasi sistem lama, dan regulasi sektor
keuangan.

\begin{table}[htbp]
\centering
\scriptsize
\caption{Contoh pemetaan visi arsitektur Bank Wanua.}
\label{tab:bank-wanua-vision}
\begin{tabularx}{\textwidth}{@{}p{0.24\textwidth} X@{}}
\hline
\textbf{Elemen} & \textbf{Rumusan awal} \\
\hline
Masalah utama & Data nasabah, proses produk, dan kanal digital belum
terintegrasi secara memadai. \\
Pemangku kepentingan & Pimpinan bank, unit bisnis, unit kepatuhan, tim
teknologi, mitra layanan digital, dan nasabah. \\
Hasil yang diharapkan & Layanan baru dapat dirancang dan diluncurkan lebih
cepat dengan kendali keamanan dan kepatuhan yang jelas. \\
Ruang lingkup & Integrasi data, orkestrasi proses, tata kelola API, dan
hubungan dengan sistem lama. \\
Batasan & Tidak membahas penggantian penuh core banking dan tidak menyusun
rancangan teknis rinci sampai level kode. \\
Kebutuhan pustaka & EA, platform digital, API, tata kelola data, modernisasi
sistem lama, dan regulasi keuangan digital. \\
\hline
\end{tabularx}
\end{table}

\section{Langkah Analisis dan Artefak Pendukung}

Bagian ini menjadi panduan kerja individual. Setiap langkah sebaiknya
menghasilkan catatan yang dapat dipakai kembali dalam makalah.

\subsection{Langkah 1: Kunci Pertanyaan Riset}

Ambil pertanyaan riset dari Bab~\ref{ch:strategi-bisnis}. Pastikan pertanyaan
tersebut tidak terlalu luas. Pertanyaan seperti ``Bagaimana transformasi
digital memengaruhi perusahaan?'' masih terlalu umum. Pertanyaan yang lebih
siap dianalisis misalnya: ``Bagaimana TOGAF-lite dapat digunakan untuk
merancang integrasi layanan digital pada bank daerah dengan mempertimbangkan
tata kelola API dan kepatuhan data?''

\subsection{Langkah 2: Peta Pemangku Kepentingan dan Concern}

Daftarkan pemangku kepentingan yang terdampak oleh perubahan arsitektur. Untuk
masing-masing pihak, tuliskan kebutuhan atau kekhawatirannya. Jangan hanya
menulis jabatan. Tuliskan kepentingan yang perlu dijawab dalam makalah.
Pimpinan mungkin peduli pada biaya dan manfaat. Unit kepatuhan peduli pada
risiko regulasi. Tim teknologi peduli pada kompleksitas integrasi. Pengguna
peduli pada kenyamanan dan keandalan layanan.

\subsection{Langkah 3: Tetapkan Prinsip Arsitektur}

Prinsip arsitektur adalah aturan berpikir yang membantu memilih solusi.
Contohnya: data penting harus memiliki pemilik yang jelas; API yang dibuka ke
mitra harus memiliki mekanisme pengendalian akses; perubahan sistem lama harus
dilakukan bertahap; dan setiap rekomendasi harus dapat dikaitkan dengan nilai
bisnis. Prinsip tidak perlu banyak. Tiga sampai lima prinsip yang tajam lebih
berguna daripada daftar panjang yang tidak dipakai.

\subsection{Langkah 4: Rumuskan Visi Arsitektur}

Tuliskan visi arsitektur dalam satu paragraf utama. Paragraf itu harus
menjawab masalah, arah perubahan, nilai yang diharapkan, ruang lingkup, dan
batasan. Jika paragraf terasa terlalu panjang, pecah menjadi pernyataan visi
dan daftar batasan. Visi yang baik tidak menjanjikan segalanya; ia memberi
arah yang cukup jelas untuk dianalisis.

\subsection{Langkah 5: Susun Strategi Pencarian Pustaka}

Strategi pencarian pustaka sebaiknya ditulis sebelum mahasiswa mengumpulkan
artikel. Mulailah dari konsep utama dalam pertanyaan riset. Untuk topik Bank
Wanua, kata kunci dapat mencakup \textit{enterprise architecture},
\textit{digital platform}, \textit{API governance},
\textit{legacy system modernization}, \textit{banking digital transformation},
\textit{data governance}, dan \textit{TOGAF}. Padankan kata kunci global
dengan istilah Indonesia bila diperlukan, misalnya ``arsitektur enterprise'',
``tata kelola data'', dan ``layanan keuangan digital''.

\subsection{Langkah 6: Seleksi dan Catat Referensi}

Setiap referensi perlu diseleksi berdasarkan relevansi, mutu, dan fungsi dalam
makalah. Buku kerangka kerja dapat dipakai untuk definisi dan konsep dasar.
Artikel jurnal membantu membangun temuan akademik. Regulasi memberi konteks
batas. Standar profesional memberi bahasa praktik. Catatan referensi minimal
memuat: tujuan sumber, konsep penting, metode atau konteks, temuan utama, dan
alasan sumber tersebut dipakai.

\subsection{Langkah 7: Buat Matriks Sintesis}

Matriks sintesis membantu mahasiswa membandingkan sumber, bukan sekadar
meringkas satu per satu. Gunakan kolom yang sesuai dengan topik. Untuk topik
platform digital, kolom dapat mencakup definisi platform, peran ekosistem,
mekanisme tata kelola, implikasi arsitektur, dan celah riset. Literatur
tentang platform dan ekosistem digital, misalnya, dapat membantu menjelaskan
mengapa arsitektur tidak hanya berbicara tentang aplikasi internal, tetapi
juga hubungan antara organisasi, pengguna, dan mitra
\cite{parker2016platform,cusumano2019business,evans2016matchmakers}.

\subsection{Langkah 8: Tulis Metode Awal}

Metode awal menjelaskan bagaimana makalah akan menggunakan TOGAF-lite.
Contohnya: penelitian konseptual ini menggunakan adaptasi ringan ADM untuk
menyusun visi arsitektur, memetakan kondisi saat ini dan target, melakukan
analisis kesenjangan, serta merumuskan roadmap awal. Tinjauan pustaka
digunakan untuk membangun kerangka konseptual dan memvalidasi kategori
analisis. Metode awal seperti ini belum final, tetapi sudah cukup untuk
mengarahkan penulisan.

\begin{table}[htbp]
\centering
\scriptsize
\caption{Artefak pendukung Bab 3 dan fungsinya dalam makalah.}
\label{tab:artefak-bab3}
\begin{tabularx}{\textwidth}{@{}p{0.25\textwidth} p{0.31\textwidth} X@{}}
\hline
\textbf{Artefak} & \textbf{Isi minimum} & \textbf{Dipakai untuk bagian makalah} \\
\hline
Visi arsitektur & Masalah, arah perubahan, nilai, ruang lingkup, batasan &
Pendahuluan, metode, dan pembahasan awal. \\
Peta stakeholder-concern & Pihak terdampak dan kepentingan utamanya & Latar
belakang, analisis kebutuhan, diskusi risiko. \\
Log pencarian pustaka & Kata kunci, sumber pencarian, kriteria seleksi &
Metode tinjauan pustaka. \\
Bibliografi kerja & Minimal sepuluh rujukan dengan catatan kegunaan &
Tinjauan pustaka dan daftar pustaka. \\
Matriks sintesis & Perbandingan konsep, metode, temuan, dan relevansi &
Kerangka konseptual dan pembahasan. \\
Outline metode & Penjelasan penggunaan TOGAF-lite dan ADM & Bagian metode atau
pendekatan penelitian. \\
\hline
\end{tabularx}
\end{table}

\section{Integrasi ke Makalah Publikasi}

Bab ini berkontribusi langsung pada tiga bagian naskah: tinjauan pustaka,
kerangka konseptual, dan metode. Visi arsitektur biasanya tidak ditempatkan
sebagai artefak mentah yang berdiri sendiri. Ia diterjemahkan menjadi narasi
dalam pendahuluan dan metode. Matriks sintesis tidak selalu perlu dimasukkan
penuh ke naskah, tetapi hasil sintesisnya harus tampak dalam kerangka
konseptual dan pembahasan.

\subsection{Dari Visi ke Pendahuluan}

Pendahuluan makalah dapat menggunakan visi arsitektur untuk menjelaskan
urgensi. Contoh kalimat:

\begin{quote}
Penelitian ini membahas kebutuhan arsitektur layanan digital pada bank daerah
yang menghadapi tekanan untuk mempercepat peluncuran layanan baru, menjaga
kepatuhan data, dan tetap mempertahankan keberlanjutan sistem lama.
\end{quote}

Kalimat tersebut lebih kuat daripada pernyataan umum bahwa ``transformasi
digital sangat penting''. Ia menunjukkan konteks, tekanan perubahan, dan batas
topik.

\subsection{Dari Pustaka ke Kerangka Konseptual}

Tinjauan pustaka perlu diolah menjadi hubungan antar-konsep. Misalnya, konsep
platform digital membantu menjelaskan hubungan antara penyedia layanan,
pengguna, dan mitra \cite{parker2016platform,cusumano2019business}. Konsep
API membantu menjelaskan bagaimana kapabilitas organisasi dapat dibuka dan
dikelola sebagai layanan \cite{jacobson2011apis,medjaoui2018continuous}.
Konsep \textit{operating model} dan arsitektur enterprise membantu menautkan
pilihan teknologi dengan proses dan tata kelola organisasi
\cite{ross2006enterprise,ross2019designed}. Hubungan antar-konsep itulah yang
menjadi kerangka konseptual, bukan sekadar daftar definisi.

\subsection{Dari ADM ke Metode}

Metode makalah dapat menjelaskan bahwa analisis dilakukan melalui adaptasi
ringan ADM. Tahapannya mencakup perumusan visi arsitektur, pemetaan kondisi
saat ini, perumusan kondisi target, analisis kesenjangan, dan penyusunan
roadmap. Karena makalah ini bersifat individual dan akademik, keluaran ADM
dipilih secara selektif. Artefak yang dipakai hanyalah yang memperjelas
argumen penelitian.

\section{Diskusi Kritis, Trade-off, dan Perangkap Umum}

\subsection{Trade-off Pemakaian TOGAF-lite}

TOGAF-lite berguna karena memberi struktur tanpa menuntut dokumentasi yang
terlalu berat. Namun, penyederhanaan memiliki risiko. Jika terlalu
disederhanakan, TOGAF hanya menjadi label. Jika terlalu formal, makalah
berubah menjadi dokumen prosedural yang sulit dibaca. Mahasiswa perlu menjaga
keseimbangan: cukup disiplin untuk menunjukkan tahapan analisis, tetapi cukup
ringkas agar naskah tetap fokus pada argumen.

\begin{table}[htbp]
\centering
\scriptsize
\caption{Trade-off dalam Bab 3.}
\label{tab:tradeoff-bab3}
\begin{tabularx}{\textwidth}{@{}p{0.23\textwidth} X X@{}}
\hline
\textbf{Keputusan} & \textbf{Jika terlalu ringan} & \textbf{Jika terlalu berat} \\
\hline
TOGAF-lite & Menjadi tempelan istilah tanpa metode jelas & Naskah penuh
artefak dan kehilangan argumen utama. \\
Tinjauan pustaka & Sumber sedikit dan klaim lemah & Terlalu banyak kutipan
tanpa sintesis. \\
Visi arsitektur & Arah perubahan kabur & Visi terlalu luas dan tidak selesai
dianalisis. \\
Konteks Indonesia & Makalah terasa generik & Pembahasan regulasi menutupi
fokus arsitektur. \\
Metode & Pembaca tidak tahu cara analisis dilakukan & Metode terlalu teknis
untuk target naskah. \\
\hline
\end{tabularx}
\end{table}

\subsection{Perangkap Umum}

Perangkap pertama adalah memakai TOGAF sebagai daftar fase yang ditempelkan ke
makalah tanpa hubungan yang nyata dengan pertanyaan riset. Hindari menulis
bahwa penelitian ``menggunakan ADM'' tetapi tidak menunjukkan bagaimana setiap
tahap membantu analisis. Perangkap kedua adalah mengumpulkan referensi hanya
karena judulnya mirip. Referensi harus dipilih karena memberi konsep, metode,
temuan, atau konteks yang benar-benar diperlukan.

Perangkap ketiga adalah menulis visi arsitektur terlalu normatif, misalnya
``membangun sistem digital yang aman, cepat, efisien, dan terintegrasi''.
Pernyataan seperti itu terdengar baik, tetapi belum analitis. Apa yang
diintegrasikan? Aman terhadap risiko apa? Cepat untuk proses siapa? Efisien
dibandingkan kondisi apa? Pertanyaan semacam ini perlu dijawab agar visi
menjadi bahan analisis, bukan slogan.

\section{Latihan Terapan Individual}

Latihan pada bab ini dikerjakan secara individual dan menjadi bahan langsung
untuk makalah publikasi.

\subsection*{Latihan 3.1: Visi Arsitektur Satu Halaman}

Tulis visi arsitektur untuk topik makalah masing-masing. Gunakan struktur
berikut:

\begin{enumerate}
  \item judul topik;
  \item masalah utama dalam dua sampai tiga kalimat;
  \item pemangku kepentingan utama dan concern-nya;
  \item pernyataan visi arsitektur;
  \item ruang lingkup dan batasan;
  \item nilai yang diharapkan;
  \item risiko atau asumsi awal.
\end{enumerate}

\subsection*{Latihan 3.2: Strategi Pencarian dan Bibliografi Awal}

Susun strategi pencarian pustaka. Tuliskan kata kunci utama, kombinasi kata
kunci, sumber pencarian, dan kriteria seleksi. Kumpulkan sedikitnya sepuluh
referensi awal. Untuk setiap referensi, tuliskan satu sampai tiga kalimat
tentang kegunaannya bagi makalah.

\subsection*{Latihan 3.3: Matriks Sintesis Lima Referensi Terkuat}

Pilih lima referensi yang paling relevan. Bandingkan kelimanya dalam matriks
yang memuat konsep utama, konteks, metode atau pendekatan, temuan penting, dan
hubungan dengan pertanyaan riset. Setelah itu, tulis satu paragraf sintesis.
Paragraf tersebut harus menunjukkan hubungan antar-sumber, bukan meringkas
sumber secara terpisah.

\subsection*{Latihan 3.4: Paragraf Metode Awal}

Tulis satu paragraf yang menjelaskan bagaimana makalah menggunakan
TOGAF-lite. Paragraf ini minimal memuat alasan memilih TOGAF-lite, tahapan
analisis yang digunakan, artefak yang dihasilkan, dan peran tinjauan pustaka
dalam analisis.

\begin{table}[htbp]
\centering
\scriptsize
\caption{Rubrik latihan Bab 3.}
\label{tab:rubrik-bab3}
\begin{tabularx}{\textwidth}{@{}p{0.25\textwidth} X X@{}}
\hline
\textbf{Komponen} & \textbf{Indikator baik} & \textbf{Perlu diperbaiki bila} \\
\hline
Visi arsitektur & Masalah, arah perubahan, nilai, ruang lingkup, dan batasan
jelas & Masih berupa slogan atau terlalu luas. \\
Stakeholder-concern & Pihak terdampak dan kepentingannya dapat dibedakan &
Hanya berisi daftar jabatan. \\
Bibliografi awal & Minimal sepuluh sumber relevan dan diberi catatan kegunaan
& Sumber dipilih hanya karena mudah ditemukan. \\
Matriks sintesis & Menunjukkan hubungan antar-sumber & Berisi ringkasan
artikel satu per satu. \\
Metode awal & Menjelaskan penggunaan TOGAF-lite secara operasional & Hanya
menyebut TOGAF tanpa tahapan analisis. \\
\hline
\end{tabularx}
\end{table}

\subsection*{Refleksi Singkat}

Jawab tiga pertanyaan berikut dalam catatan pribadi:

\begin{enumerate}
  \item Bagian mana dari visi arsitektur yang masih paling kabur?
  \item Referensi apa yang paling membantu mengubah cara Anda melihat masalah?
  \item Apa yang perlu dipersempit agar makalah dapat selesai sebagai naskah
        publikasi individual?
\end{enumerate}

\section{Ringkasan, Refleksi, dan Jembatan}

\subsection*{Poin Kunci}

\begin{itemize}
  \item TOGAF-lite digunakan sebagai kerangka berpikir ringan, bukan sebagai
        kewajiban membuat seluruh artefak TOGAF.
  \item ADM membantu mahasiswa menata analisis dari masalah menuju visi,
        kondisi saat ini, kondisi target, kesenjangan, dan rekomendasi.
  \item Visi arsitektur memberi arah perubahan yang jelas dan membatasi ruang
        lingkup makalah.
  \item Tinjauan pustaka memberi dasar akademik untuk konsep, metode, dan
        argumen.
  \item Bibliografi kerja dan matriks sintesis adalah jembatan antara aktivitas
        membaca dan penulisan naskah.
\end{itemize}

\subsection*{Menjawab Pertanyaan Pemandu}

Perubahan arsitektur yang dibahas dalam makalah harus dinyatakan sebagai visi
yang spesifik, bukan keinginan umum untuk ``bertransformasi digital''.
Pemangku kepentingan perlu dipetakan karena setiap perubahan arsitektur
membawa dampak yang berbeda bagi pimpinan, unit bisnis, unit kepatuhan, tim
teknologi, mitra, dan pengguna. ADM dapat disederhanakan menjadi alur analisis
yang menjawab masalah, visi, baseline, target, kesenjangan, dan rekomendasi.
Literatur dipilih untuk memperkuat konsep, metode, konteks, dan pembahasan.
Visi arsitektur dan tinjauan pustaka menyatu ketika keduanya digunakan untuk
menulis kerangka konseptual dan metode makalah.

\subsection*{Uji Mandiri}

\begin{enumerate}
  \item Jelaskan perbedaan antara menggunakan TOGAF secara penuh dan
        menggunakan TOGAF-lite untuk makalah individual.
  \item Ambil satu pertanyaan riset Anda. Tuliskan tiga kata kunci global dan
        tiga kata kunci Indonesia yang relevan untuk pencarian pustaka.
  \item Pilih satu referensi terkuat dari bibliografi awal. Jelaskan apakah
        referensi itu dipakai untuk definisi, kerangka konsep, metode,
        konteks, atau pembahasan.
\end{enumerate}

\subsection*{Jembatan ke Bab Berikutnya}

Bab berikutnya masuk ke arsitektur bisnis. Jika Bab 3 menjawab arah perubahan
dan fondasi akademik, Bab 4 mulai membedah proses, kapabilitas, nilai, dan
pemangku kepentingan bisnis secara lebih rinci. Visi arsitektur yang sudah
dibuat pada bab ini akan dipakai untuk memilih kapabilitas mana yang penting,
proses mana yang bermasalah, dan nilai bisnis apa yang harus dijaga dalam
rekomendasi makalah.

\section{Bacaan Lanjutan dan Lampiran}

\subsection*{Bacaan Inti}

\begin{itemize}
  \item The Open Group Architecture Framework versi terbaru sebagai rujukan
        utama TOGAF dan ADM \cite{togaf10}.
  \item Ross, Weill, dan Robertson untuk memahami hubungan antara enterprise
        architecture, operating model, dan fondasi eksekusi
        \cite{ross2006enterprise}.
  \item Ross, Beath, dan Mocker untuk melihat bagaimana organisasi merancang
        ulang dirinya dalam ekonomi digital \cite{ross2019designed}.
\end{itemize}

\subsection*{Bacaan untuk Memperdalam Topik}

\begin{itemize}
  \item Untuk topik platform dan ekosistem digital, baca rujukan tentang
        platform dan pasar dua sisi
        \cite{parker2016platform,cusumano2019business,evans2016matchmakers}.
  \item Untuk topik API dan integrasi layanan, baca rujukan tentang API
        sebagai produk dan tata kelola API
        \cite{jacobson2011apis,medjaoui2018continuous}.
  \item Untuk topik modernisasi sistem dan arsitektur perangkat lunak, baca
        pembahasan tentang monolith dan microservices \cite{newman2021building}.
  \item Untuk konteks Indonesia, sesuaikan rujukan dengan topik: perlindungan
        data pribadi \cite{uupdp2022}, penyelenggaraan sistem elektronik
        \cite{pp712019}, Satu Data Indonesia \cite{satudata}, atau standar API
        pembayaran nasional \cite{snapbi}.
\end{itemize}

\subsection*{Lampiran Kerja Bab 3}

Gunakan templat ringkas berikut sebagai catatan kerja:

\begin{description}
  \item[Topik makalah] Tuliskan topik dalam satu kalimat.
  \item[Pertanyaan riset] Tuliskan pertanyaan riset utama dan, bila perlu,
        satu subpertanyaan.
  \item[Visi arsitektur] Tuliskan arah perubahan dalam satu paragraf.
  \item[Stakeholder-concern] Tuliskan minimal empat pihak dan concern
        masing-masing.
  \item[Kata kunci pustaka] Tuliskan kata kunci utama dalam bahasa Inggris dan
        Indonesia.
  \item[Bibliografi awal] Tuliskan minimal sepuluh referensi beserta catatan
        kegunaan.
  \item[Metode awal] Jelaskan bagaimana TOGAF-lite dipakai dalam analisis.
\end{description}
\setlength{\emergencystretch}{0pt}
